\documentclass[11pt]{article}
\usepackage[margin=1in]{geometry}
\usepackage{amsmath,amssymb,amsthm}
\usepackage{hyperref}
\hypersetup{colorlinks=true,urlcolor=blue,linkcolor=blue}
\newtheorem{theorem}{Theorem}
\newtheorem{lemma}{Lemma}
\newcommand{\C}{\mathbb C}
\newcommand{\bc}[1]{\overline{#1}}
\title{The concentric ellipse of common Simson points\\
\large MathDB 352211: a full proof}
\author{Principia Math In Progress}
\date{21 September 2026}
\begin{document}
\maketitle
\section*{Task 1: statement and scope}
Let a nondegenerate Poncelet triangle family be inscribed in a circle and
circumscribed about a fixed nested ellipse. For each fixed point $F$ on the
circle, consider the Simson lines of $F$ relative to the triangles.
Conjecture 1 of \cite{RG}, or Conjecture 2.1 in its journal version, asks
whether their envelope point traces an ellipse concentric with the caustic
as $F$ traverses the circle. We prove this and give its semiaxes.

\begin{theorem}
Normalize the outer circle to $|z|=1$, and denote the caustic's foci by
$f,g\in\C$. For $|u|=1$, the Simson lines of $F=u$ form the full pencil
through
\[
             W(u)=\frac{u+f+g-fg/u}{2}.
\]
As $u$ traverses the circle, $W(u)$ traces the complete nondegenerate ellipse
with center $(f+g)/2$ and semiaxes
$\bigl(1+|fg|\bigr)/2$ and $\bigl(1-|fg|\bigr)/2$.
A circle is included as an ellipse.
\end{theorem}

\section*{1. The Simson line in complex coordinates}
Write $s=f+g$, $p=fg$. For a triangle with vertices $a,b,c$ on the unit
circle put $\lambda=abc$, $s_1=a+b+c$, $s_2=ab+ac+bc$.
We use the standard isogonal-foci identities
\begin{equation}\label{foci}
       s_1=s+\bc p\lambda,\qquad s_2=p+\bc s\lambda.
\end{equation}
These are the input used in the proof of Theorem 1 of \cite{RG}
(Theorem 2.1 in its journal version), and are displayed as equation (3.1)
in \cite{GG}. They follow because the two foci of an inellipse are
isogonal conjugates: in unit-circle coordinates their relation is
$f+g+\lambda\bc f\bc g=a+b+c$. Conjugation and multiplication by $\lambda$
give the second identity. We claim no new result for these identities.

The line through $a,b$ has equation $z+ab\bc z=a+b$. Its perpendicular
foot from $u$ is
\[
 q_{ab}=\frac{u+a+b-ab/u}{2};
\]
indeed it lies on that line and
$u-q_{ab}=ab\,\bc{(u-q_{ab})}$, the normal-direction condition.
Substitution, using $\bc a=1/a$ and cyclically, shows that all three
feet lie on
\begin{equation}\label{simson}
 z-\frac{\lambda}{u}\bc z
       =\frac{u+s_1-s_2/u-\lambda/u^2}{2}.
\end{equation}
At least two feet are distinct: otherwise all three perpendiculars would
meet the three nonparallel sidelines at the same point, which is
impossible for a nondegenerate triangle. Thus \eqref{simson} is the Simson
line, including when $u$ is a vertex.

\newpage
\section*{2. Its common point and the full pencil}
Set $W=(u+s-p/u)/2$. Because $|u|=1$,
$\bc W=(1/u+\bc s-\bc p u)/2$.
Substitution of \eqref{foci} into \eqref{simson} gives exactly
\begin{equation}\label{pencil}
            z-W=\frac{\lambda}{u}\bc{(z-W)}.
\end{equation}
In particular the point $W$ is independent of the triangle.

The parameter $\lambda$ traverses the whole unit circle. To see this
directly, the vertex $a$ is a root of
\[
 z^3-(s+\bc p\lambda)z^2+(p+\bc s\lambda)z-\lambda=0.
\]
Consequently
\[
 \lambda=B(a):=\frac{a(a-f)(a-g)}
                       {(1-\bc f a)(1-\bc g a)}.
\]
Here $|f|,|g|<1$ because the caustic is nested in the circle. Each factor
$(a-h)/(1-\bc h a)$ maps the unit circle onto itself with winding number
one, as does $a$. Hence $B$ has winding number three and is onto. The
Poncelet family lets $a$ traverse the outer circle, proving the assertion.
In fact the derivative of the argument of $B(e^{it})$ is
\[
 1+\frac{1-|f|^2}{|e^{it}-f|^2}
   +\frac{1-|g|^2}{|e^{it}-g|^2}>0.
\]
Thus \eqref{pencil} describes every unoriented line through $W$. Its
envelope is the single point $W$; equivalently the equations of a line
and its derivative with respect to its angle have the unique solution
$W$.

For completeness this point is also the antipode of $u$ on the circular
locus of the perpendicular projection $V$ of $u$ onto the Simson line.
Reflection in \eqref{pencil} sends $z$ to
$W+(\lambda/u)(\bc z-\bc W)$, so
\[
 V=\frac{u+W}{2}+\frac{\lambda}{2u}(\bc u-\bc W).
\]
This is the full circle of center $(u+W)/2$ and radius $|u-W|/2$.
It is nondegenerate, since
\[
                    u-W=\frac{(u-f)(u-g)}{2u}\ne0.
\]
The source already proves that this projection locus is a circle. Our
calculation identifies the common point needed for the ellipse theorem.

\section*{3. The complete concentric ellipse}
As $u$ runs through $|u|=1$,
\begin{equation}\label{ellipse}
             W-\frac s2=\frac{u-p\bc u}{2}.
\end{equation}
Writing $p=p_r+ip_i$, the real linear map on the right has matrix
\[
 M=\frac12
 \begin{pmatrix}1-p_r&-p_i\\-p_i&1+p_r\end{pmatrix}.
\]
Its determinant is $(1-|p|^2)/4>0$, and its positive eigenvalues are
$(1-|p|)/2$ and $(1+|p|)/2$. The image of the full unit circle under
this invertible map is precisely the full ellipse of these semiaxes.
Its center is $s/2$, the midpoint of the caustic's two foci, hence the
caustic's center. This proves every assertion of the theorem.

\newpage
\section*{4. Explicit equation, scale, and verification}
An implicit equation without choices of axis is obtained by putting
$v=2(W-s/2)$. Equation \eqref{ellipse} implies
\[
 v+p\bc v=(1-|p|^2)u,\qquad
          |v+p\bc v|^2=(1-|p|^2)^2.
\]
Conversely the real map $v\mapsto v+p\bc v$ is invertible because
$|p|<1$, so this equation introduces no extra real component.
If $p=\rho e^{i\phi}\ne0$, rotate coordinates by $\phi/2$ and write
$u=e^{i\phi/2}e^{it}$. Then
\[
 e^{-i\phi/2}(W-s/2)
     =\frac{1-\rho}{2}\cos t
        +i\frac{1+\rho}{2}\sin t.
\]
For $p=0$ the locus is the circle of radius $1/2$.

Restoring an outer circle of center $O$ and radius $R$, with caustic
foci $F_1,F_2$, the center of the locus is $(F_1+F_2)/2$ and its
semiaxes are
\[
             \frac{R^2+|F_1-O|\,|F_2-O|}{2R},
       \qquad
             \frac{R^2-|F_1-O|\,|F_2-O|}{2R}.
\]
The minor semiaxis is strictly positive. The assertion requires neither
concentric original conics nor alignment of their axes.

The accompanying \texttt{verify.py} checks the exact identities for
the feet, the Simson line, its common point, and the projection circle.
It also checks the matrix and its inverse relation symbolically.
Independent Euclidean projections check 150 configurations from five
Poncelet families at 90 decimal digits. These include circular caustics
and a focus at the outer center. The finite checks supplement the
general proof.

The common-point assertion recorded separately as MathDB 352210 is a
lemma here and receives no separate solution credit. Classical and
bicentric Simson concurrence precedents are acknowledged. The count is
one for the complete concentric-ellipse target.

\section*{Proposed follow-up Tasks 2--5}
\begin{enumerate}
\item \textbf{Task 2.} Characterize exactly when this ellipse is
axis-aligned with the caustic, using its two foci and the circumcenter.
\item \textbf{Task 3.} Classify every tangency and intersection between
the Simson-point ellipse and the caustic.
\item \textbf{Task 4.} Determine which concentric ellipse pairs can occur
as a caustic and its Simson-point locus.
\item \textbf{Task 5.} Determine the envelopes when the projection
direction is obtained by a fixed rotation of each side's normal.
\end{enumerate}

\begin{thebibliography}{9}
\bibitem{RG} D. Reznik and R. Garcia,
\emph{Poncelet Parabola Pirouettes}.
International Journal of Geometry \textbf{12}(1) (2023), 93--116.
\href{https://arxiv.org/abs/2110.06356}{arXiv:2110.06356v5}.
Both complete 24-page versions were read; the target remains conjectural.
\bibitem{GG} R. A. Garcia and L. G. Gheorghe,
\emph{Poncelet Pairs via Isogonal Conjugacy}.
Aequationes Mathematicae \textbf{100}, article 31 (2026).
\href{https://doi.org/10.1007/s00010-026-01274-w}{doi:10.1007/s00010-026-01274-w}.
Equation (3.1) records the standard foci identities.
\end{thebibliography}
\end{document}
